\section{Introduction}
\label{sec:introduction}

The JSON Canonicalization Scheme (\rfcjcs)~\cite{rfc8785} standardizes a
deterministic serialization of JSON values for use in digital signatures,
content-addressable storage, and distributed consensus protocols.  A
critical requirement is that floating-point numbers must be formatted
identically to the output of \ecma \texttt{Number::toString}~\cite{ecma262},
which specifies shortest-representation formatting under roundTiesToEven
semantics.

Several deployment contexts place canonicalization on the critical path.
In high-frequency signature verification, every document must be
canonicalized before its signature can be checked, making formatting
throughput a direct bottleneck.  Content-addressed storage systems hash
the canonical form at both ingestion and retrieval, so canonicalization
cost scales linearly with object throughput.  Distributed consensus
protocols require byte-identical serialization for state-transition
agreement, invoking canonicalization on every proposal and vote.  In
each scenario, a constant-factor improvement in the formatting stage
translates directly to throughput gains: the 1.4--1.8$\times$ speedup
we measure on number-dense payloads corresponds to freeing
approximately 24\% of a CPU core at 1{,}000 canonicalizations per
second on a representative 14\,KB workload.

Beyond throughput, the choice of digit-generation algorithm is a
correctness concern in systems where canonicalized bytes are hashed and
signed at every layer of a processing pipeline.  In deterministic
infrastructure---where data-processing steps execute in sandboxed
isolation and every governed artifact is schema-validated, canonically
serialized, and cryptographically bound---a formatting divergence at the
number-rendering layer does not produce one wrong document: it
invalidates every hash computed over those bytes, every signature
computed over those hashes, and every audit trail that chains those
signatures.  The digit-generation algorithm sits at the foundation of a
determinism guarantee that no higher layer can independently recover.
This makes both the correctness validation (conformance gating) and the
performance characterization reported here directly relevant to
infrastructure where canonicalization is not optional but
load-bearing~\cite{rfc8785}.

Two families of algorithms can satisfy this requirement: multiprecision
arithmetic approaches exemplified by \burgerdybvig~\cite{burger1996printing},
and fixed-width integer approaches exemplified by
\schubfach~\cite{giulietti2022schubfach},
\ryualg~\cite{adams2018ryu}, and \dragonbox~\cite{jeon2021dragonbox}.
While these algorithms are individually well-studied, no prior work has
empirically compared them \emph{inside a specification-constrained
canonicalization pipeline} where conformance is a hard prerequisite and
the formatting stage is only one component of the total cost.

This paper makes three contributions:

\begin{enumerate}
\item A benchmark methodology that gates performance measurement on
      conformance verification: implementations must pass 286{,}362
      oracle vectors and 776 structural conformance cases before timing
      data is collected.  We document noise floors and minimum observable
      effect sizes for all comparisons.

\item The first empirical comparison of four digit-generation algorithms
      (three fixed-width: \schubfach, \ryualg, \dragonbox; one
      multiprecision: \burgerdybvig) under \ecma/\rfcjcs constraints,
      showing that all three fixed-width algorithms perform within
      1--8\% of each other while each achieves 1.4--1.8$\times$ speedup
      over \burgerdybvig on number-dense workloads ($\nu(D) \geq 40\%$)
      in Go API benchmarks; the advantage replicates at
      2.3--3.0$\times$ in Rust.
      String-dominant and deeply-nested inputs show $<$5\% difference.

\item A cross-language validation design (4$\times$2 matrix of
      algorithm $\times$ language) that separates algorithmic effects from
      language/runtime effects, with Go cross-compilation determinism
      verified for x86-64 and arm64 via QEMU user-mode emulation.
      All evidence is released as a reproducibility artifact.
\end{enumerate}
